\section{Surrounded Regions}
\label{sec:graph-surrounded-regions}

\problemstatement{Given an $m \times n$ board of \texttt{'X'} and
\texttt{'O'}, \textbf{capture} every region of \texttt{'O'}s that is
fully \textbf{surrounded} by \texttt{'X'}s -- flip every captured
\texttt{'O'} to \texttt{'X'}. A region is \textbf{not} captured if any
\texttt{'O'} in it touches the border of the board.}

\begin{patternbox}
\emph{``Mark everything reachable from the boundary first, then
process by elimination''} -- rather than trying to directly test
whether each interior region is ``surrounded'' (which would require
checking every cell of every region against the border), it is far
simpler to run a traversal \textbf{only from the border's
\texttt{'O'}s}: anything that traversal reaches is, by definition,
connected to the border and therefore \emph{safe}. Everything left
unmarked afterward is unreachable from the border -- the surrounded
regions -- and gets flipped.
\end{patternbox}

\subsection*{Algorithm}

\begin{enumerate}
  \item For every \texttt{'O'} on the board's \textbf{border} (first
        row, last row, first column, last column), run DFS (or BFS),
        marking every $4$-directionally connected \texttt{'O'} with a
        temporary marker (e.g.\ \texttt{'\#'}) instead of visiting a
        separate \textit{visited} array.
  \item Scan the entire board: any remaining \texttt{'O'} (never
        reached from the border) is surrounded -- flip it to
        \texttt{'X'}. Any \texttt{'\#'} (reached from the border, safe)
        gets restored to \texttt{'O'}.
\end{enumerate}

\subsection*{C++ implementation}

\begin{lstlisting}
#include <bits/stdc++.h>
using namespace std;

void markSafe(int r, int c, vector<vector<char>>& board) {
    int rows = board.size(), cols = board[0].size();
    if (r < 0 || r >= rows || c < 0 || c >= cols) return;
    if (board[r][c] != 'O') return;

    board[r][c] = '#';
    markSafe(r - 1, c, board);
    markSafe(r + 1, c, board);
    markSafe(r, c - 1, board);
    markSafe(r, c + 1, board);
}

void solve(vector<vector<char>>& board) {
    int rows = board.size(), cols = board[0].size();

    for (int r = 0; r < rows; r++) {
        markSafe(r, 0, board);
        markSafe(r, cols - 1, board);
    }
    for (int c = 0; c < cols; c++) {
        markSafe(0, c, board);
        markSafe(rows - 1, c, board);
    }

    for (int r = 0; r < rows; r++) {
        for (int c = 0; c < cols; c++) {
            if (board[r][c] == 'O') board[r][c] = 'X';
            else if (board[r][c] == '#') board[r][c] = 'O';
        }
    }
}

int main() {
    vector<vector<char>> board = {
        {'X', 'X', 'X', 'X'},
        {'X', 'O', 'O', 'X'},
        {'X', 'X', 'O', 'X'},
        {'X', 'O', 'X', 'X'}
    };
    solve(board);
    for (auto& row : board) {
        for (char x : row) cout << x << " ";
        cout << endl;
    }
    return 0;
}
\end{lstlisting}

\subsection*{Dry run}

\dryrun{Take the $4\times4$ board above.}

Border \texttt{'O'}s: only $(3,1)$ is on the border (last row) -- but
checking the full border scan, none of row $0$, row $3$ (except
$(3,1)$, which is \texttt{'O'}), column $0$, or column $3$ contain
other \texttt{'O'}s. $\textit{markSafe}(3,1)$ marks just $(3,1)$ as
\texttt{'\#'} (its neighbors $(2,1)=\text{X}$, $(3,0)=\text{X}$,
$(3,2)=\text{X}$ are all \texttt{'X'}, so the safe region is just this
one cell). The interior region $\{(1,1),(1,2),(2,2)\}$ is never
touched by any border-seeded traversal, since it doesn't connect to
$(3,1)$ (blocked by \texttt{'X'} on all sides). Final scan: the
interior \texttt{'O'}s flip to \texttt{'X'}; $(3,1)$'s \texttt{'\#'}
reverts to \texttt{'O'}.

\begin{complexitybox}
\textbf{Time Complexity.} $O(\textit{rows} \times \textit{cols})$ --
each cell is visited a constant number of times (once by the border
traversal if reachable, plus the one final scan).
\textbf{Space Complexity.} $O(\textit{rows} \times \textit{cols})$
recursion stack in the worst case.
\end{complexitybox}

\begin{takeawaybox}
``Process by elimination from the boundary'' is a reusable trick
whenever a problem's condition is phrased in terms of \emph{not}
reaching some special set of cells (the border, here) -- traverse
\emph{from} that special set first, and treat everything unreached as
the answer. The next section, Number of Enclaves, reuses this exact
same boundary-first elimination on a near-identical setup.
\end{takeawaybox}
